\documentclass[12pt]{article}

\usepackage{amsmath, amssymb, amsthm, enumerate, graphicx}
\usepackage[usenames,dvipsnames]{color}
\usepackage{bm}
\usepackage[colorlinks=true,urlcolor=blue]{hyperref}
\usepackage{geometry}
\geometry{margin=1in}
\usepackage{float}
\usepackage{graphics}
\setlength{\marginparwidth}{2.15cm}
\usepackage{booktabs}
\usepackage{enumitem}
\usepackage{epsfig}
\usepackage{setspace}
\usepackage{parskip}
\usepackage[normalem]{ulem}
\usepackage{tikz}
\usetikzlibrary{positioning, arrows, automata}
\usepackage{pgfplots}
\usepackage[font=scriptsize]{subcaption}
\usepackage[]{algorithm2e}
\usepackage{environ}
\usepackage{bbm}
\usepackage{titling}
\usepackage{url}
\usepackage{xcolor}
\usepackage{lipsum}
\usepackage{lastpage}
\usepackage{multicol}
\usepackage{tabularx}
\usepackage{comment}
\usepackage[utf8]{inputenc}
\usepackage{setspace}
\usepackage{marvosym}
\usepackage{wrapfig}
\usepackage{datetime}
\usepackage[many]{tcolorbox}
\usepackage{array}
\usepackage{multirow}
\usepackage{wasysym}
\usepackage{cancel}
\usepackage{listings}
\usepackage[thinlines]{easytable}

\newcommand{\R}{\mathbb{R}}
\newcommand{\blackcircle}{\tikz\draw[black,fill=black] (0,0) circle (1ex);}
\renewcommand{\circle}{\tikz\draw[black] (0,0) circle (1ex);}

\usetikzlibrary{positioning,calc}

\newtcolorbox[]{solution}[1][]{%
    breakable,
    enhanced,
    colback=white,
    %title=Solution,
    #1
}

\begin{document}
\section*{}
\begin{center}
  \centerline{\textsc{\LARGE  Homework 6 Submission Template}}
  \vspace{0.5em}
  \textsc{\large CMU 10-703: Deep Reinforcement Learning (Fall 2026)}
  \vspace{0.5em}
  \centerline{Humanoid Control (Motion Imitation) with Deep RL}
\end{center}

Use this template to record your answers for Homework 6. Add your answers using \LaTeX and then save your document as a PDF to upload to Gradescope. You are required to use this template to submit your answers. \textbf{You should not alter this template in any way} other than to insert your solutions. You must submit all \pageref{LastPage} pages of this template to Gradescope. Do not remove the instructions page(s). Altering this template or including your solutions outside of the provided boxes can result in your assignment being graded incorrectly.

Also submit your executed Jupyter notebook with all cells run and outputs visible.

\newpage

\section*{Problem 0: Collaborators}
Enter your team's names and Andrew IDs in the boxes below. If you do not do this, you may lose points on your assignment.

Name 1: \begin{tcolorbox}[fit,height=1cm, width=5cm, blank, borderline={1pt}{1pt},nobeforeafter]
    \begin{center}
    \vspace{3mm}
    \large{Arnav Goyal}
    \end{center}
\end{tcolorbox}
Andrew ID 1: \begin{tcolorbox}[fit,height=1cm, width=5cm, blank, borderline={1pt}{1pt},nobeforeafter]
    \begin{center}
    \vspace{3mm}
    \large{arnavgoy}
    \end{center}
\end{tcolorbox}
    \\
Name 2: \begin{tcolorbox}[fit,height=1cm, width=5cm, blank, borderline={1pt}{1pt},nobeforeafter]
    \begin{center}
    \vspace{3mm}
    \large{Vineeth Bhat}
    \end{center}
\end{tcolorbox}
Andrew ID 2: \begin{tcolorbox}[fit,height=1cm, width=5cm, blank, borderline={1pt}{1pt},nobeforeafter]
    \begin{center}
    \vspace{3mm}
    \large{vineethb}
    \end{center}
\end{tcolorbox} \\
Name 3: \begin{tcolorbox}[fit,height=1cm, width=5cm, blank, borderline={1pt}{1pt},nobeforeafter]
    \begin{center}
    \vspace{3mm}
    \large{Kushal Kejriwal}
    \end{center}
\end{tcolorbox}
Andrew ID 3: \begin{tcolorbox}[fit,height=1cm, width=5cm, blank, borderline={1pt}{1pt},nobeforeafter]
    \begin{center}
    \vspace{3mm}
    \large{kkejriwa}
    \end{center}
\end{tcolorbox} \\
\vspace{0.5cm}

\newpage

%% ====================================================================
%%  Part 2.2: Training
%% ====================================================================
\section*{Part 2.2: Train the Model}

Paste the training curves (mean reward and mean episode length) from the notebook.
\begin{solution}[height=10cm]
% YOUR SOLUTION: paste training curves here
\end{solution}


\newpage
%% ====================================================================
%%  Part 2.3: Conceptual Questions
%% ====================================================================
\section*{Part 2.3: Conceptual Questions}

\subsection*{Question 1 (Observation Space)}

\subsubsection*{1.1: What would happen if base linear velocity were included in the actor but unavailable at deployment?}
\begin{solution}[height=7cm]
% YOUR SOLUTION
\end{solution}

\subsubsection*{1.2: How does observation noise during training help with sim-to-sim (or sim-to-real) transfer?}
\begin{solution}[height=7cm]
% YOUR SOLUTION
\end{solution}

\subsubsection*{1.3: The actor does not observe base linear velocity, but the critic does. Explain why this asymmetric design helps transfer.}
\begin{solution}[height=7cm]
% YOUR SOLUTION
\end{solution}

\subsection*{Question 2 (Terminal Conditions)}

\subsubsection*{2.1: Why do we use $z$-only position thresholds instead of full 3D distance?}
\begin{solution}[height=7cm]
% YOUR SOLUTION
\end{solution}

\subsubsection*{2.2: What would go wrong if we used a tight 3D threshold early in training?}
\begin{solution}[height=7cm]
% YOUR SOLUTION
\end{solution}

\subsubsection*{2.3: What will happen if we remove all the terminal conditions? Discuss how removing terminations affects sample efficiency.}
\begin{solution}[height=7cm]
% YOUR SOLUTION
\end{solution}

\subsection*{Question 3 (Domain Randomization)}

\subsubsection*{3.1: What parameters in the deployment simulator might differ from the training simulator?}
\begin{solution}[height=7cm]
% YOUR SOLUTION
\end{solution}

\subsubsection*{3.2: What would happen if you randomized friction too aggressively (e.g., range 0.01--10.0)?}
\begin{solution}[height=7cm]
% YOUR SOLUTION
\end{solution}


\newpage
%% ====================================================================
%%  Part 3: Sim-to-Sim Deployment
%% ====================================================================
\section*{Part 3: Sim-to-Sim Deployment}

\subsection*{Question 1 (Deploy in the Same Simulator):}
Paste the screenshot from the notebook.
\begin{solution}[height=10cm]
% YOUR SOLUTION: paste screenshot here
\end{solution}

\subsection*{Question 2 (Deploy in Vanilla MuJoCo --- Sim-to-Sim Transfer):}
Paste the screenshot from the notebook.
\begin{solution}[height=10cm]
% YOUR SOLUTION: paste screenshot here
\end{solution}

\newpage

\subsection*{Question 3 (Quantitative Evaluation):}
Paste the tracking error plots from the notebook (Section~3.4) and answer:

\subsubsection*{3.1: Report metrics for both Warp and vanilla MuJoCo (anchor pos/ori error, body pos/ori error, survival rate).}
\begin{solution}[height=10cm]
% YOUR SOLUTION: paste plots and report metrics
\end{solution}

\subsubsection*{3.2: How do the Warp and vanilla MuJoCo tracking errors compare? Where is the sim-to-sim gap largest?}
\begin{solution}[height=7cm]
% YOUR SOLUTION
\end{solution}

\subsection*{Question 4 (Detailed Metrics):}
Paste the detailed metric plots from the notebook (Section~3.5) and answer:

\subsubsection*{4.1: Report MPKPE, R-MPKPE, EE errors, and joint velocity error for both Warp and vanilla MuJoCo.}
\begin{solution}[height=10cm]
% YOUR SOLUTION: paste plots and report metrics
\end{solution}

\subsubsection*{4.2: What is the difference between MPKPE and R-MPKPE? Why is R-MPKPE useful as a separate metric?}
\begin{solution}[height=7cm]
% YOUR SOLUTION
\end{solution}


\newpage
%% ====================================================================
%%  Part 4: Ablation Experiments
%% ====================================================================
\section*{Part 4: Ablation Experiments}

\subsection*{Question 1 (No Observation Noise):}

\subsubsection*{1.1: Paste the training curves for this ablation.}
\begin{solution}[height=10cm]
% YOUR SOLUTION: paste training curves here
\end{solution}

\subsubsection*{1.2: Report the survival rate (\%) and mean body position error (m).}
\begin{solution}[height=4cm]
% YOUR SOLUTION
\end{solution}

\subsubsection*{1.3: Did removing observation noise help or hurt training performance? What about deployment performance? Explain.}
\begin{solution}[height=7cm]
% YOUR SOLUTION
\end{solution}

\newpage

\subsection*{Question 2 (No Domain Randomization):}

\subsubsection*{2.1: Paste the training curves for this ablation.}
\begin{solution}[height=10cm]
% YOUR SOLUTION: paste training curves here
\end{solution}

\subsubsection*{2.2: Report the survival rate (\%) and mean body position error (m).}
\begin{solution}[height=4cm]
% YOUR SOLUTION
\end{solution}

\subsubsection*{2.3: How does the no-DR policy perform at deployment compared to the baseline? Explain the gap (if any).}
\begin{solution}[height=7cm]
% YOUR SOLUTION
\end{solution}

\newpage

\subsection*{Question 3 (No Early Termination):}

\subsubsection*{3.1: Paste the training curves for this ablation.}
\begin{solution}[height=10cm]
% YOUR SOLUTION: paste training curves here
\end{solution}

\subsubsection*{3.2: Report the survival rate (\%) and mean body position error (m).}
\begin{solution}[height=4cm]
% YOUR SOLUTION
\end{solution}

\subsubsection*{3.3: Why does the policy still learn to track (or fail to) without early termination? Relate your answer to what you discussed in 2.3.}
\begin{solution}[height=7cm]
% YOUR SOLUTION
\end{solution}


\newpage
%% ====================================================================
%%  Part 5: Conceptual Problems Part 2
%% ====================================================================
\section*{Part 5: More Conceptual Problems (0 pts)}

\subsection*{Question 4 (Reward Sign and Termination):}

\subsubsection*{4.1}
\begin{solution}[height=6cm]
% YOUR SOLUTION
\end{solution}

\subsubsection*{4.2}
\begin{solution}[height=6cm]
% YOUR SOLUTION
\end{solution}

\subsubsection*{4.3}
\begin{solution}[height=6cm]
% YOUR SOLUTION
\end{solution}

\subsection*{Question 5 (Observability and Drift):}

\subsubsection*{5.1}
\begin{solution}[height=6cm]
% YOUR SOLUTION
\end{solution}

\subsubsection*{5.2}
\begin{solution}[height=6cm]
% YOUR SOLUTION
\end{solution}

\subsubsection*{5.3}
\begin{solution}[height=6cm]
% YOUR SOLUTION
\end{solution}

\newpage
%% ====================================================================
%%  Feedback
%% ====================================================================
\section*{Feedback}

\textbf{Feedback}: You can help the course staff improve the course for future semesters by providing feedback. What was the most confusing part of this homework, and what would have made it less confusing?
\begin{solution}[height=4cm]
\end{solution}

\textbf{Collaboration}: Detail the work division amongst your group below.
\begin{solution}[height=4cm]
\end{solution}

\noindent\textbf{Time Spent}: How many hours did you spend working on this assignment? Your answer will not affect your grade. Please average your answer over all the members of your team.
\begin{solution}[height=4cm]
\begin{table}[H]
    \centering
    \begin{tabular}{r|c}
        Alone &  \hspace{3em} %ANSWER HERE%
        \\ \hline
        With teammates & \hspace{3em} %ANSWER HERE%
        \\ \hline
        With other classmates & \hspace{3em} %ANSWER HERE%
        \\ \hline
        At office hours & \hspace{3em} %ANSWER HERE%
        \\ \hline
    \end{tabular}
\end{table}
\end{solution}


\end{document}
